\subsection{The Active--Passive Boundary}

The central finding of this paper is a clean partition: KV-cache geometry detects \emph{active} cognitive states involving output suppression (deception, refusal, withholding) but not \emph{passive} states (confabulation). The unifying mechanism is \emph{output suppression}: whenever the model withholds its full representational capacity---whether because it is declining a harmful request, declining an impossible request, or constructing a deceptive response that suppresses the richer honest representation---the cache geometry becomes sparser per token. This partition maps directly onto a mechanistic distinction: active suppression creates measurable geometric perturbations, while passive states (where the model simply lacks information) produce only a faint geometric echo that cannot be reliably exploited for detection.

This has immediate consequences for AI safety monitoring. Organizations deploying internal state monitoring can detect \emph{output suppression}---cases where a model withholds its full representational capacity, whether through refusal, deception, or any form of constrained output---but cannot detect \emph{ignorance}---cases where a model is confidently wrong. Confabulation detection fundamentally requires external verification (retrieval-augmented generation, citation checking, knowledge graph alignment), not internal monitoring.

\subsection{What ``Honest Processing Is Richer'' Means}

Four independent lines of evidence (\Cref{sec:honest_richer}) converge on the finding that honest response generation produces geometrically richer KV-cache representations than deceptive or suppressive processing. This is initially counterintuitive: deception is often assumed to require \emph{more} cognitive effort (represent truth + construct lie), so one might expect richer geometry under deception.

The resolution lies in the distinction between \emph{internal effort} and \emph{representational output}. Deception may involve more internal computation, but the \emph{output representations}---the key-value pairs written to the cache---are constrained. The model must suppress the richer honest representation and substitute a narrower deceptive one. The cache captures what the model \emph{commits to}, not what it \emph{considers}. Honest commitment is richer because the model can use its full representational capacity; deceptive commitment is sparser because the model must restrict its output to maintain the deception.

This interpretation predicts that any state involving output suppression---declining to use available capacity---should produce sparser geometry. The refusal and impossibility results confirm this: both produce reduced angular spread and lower norm-per-token regardless of whether the suppression reflects safety concerns (harmful refusal) or logical impossibility (impossibility refusal). Abliterated models processing harmful content show residual suppression geometry despite having the refusal direction removed, suggesting that suppression is not solely mediated by the refusal direction but involves broader output-constraining circuitry. The common thread is restraint, not content type.

\subsection{The Single Misalignment Axis}

The discovery that deception, sycophancy, and confabulation all perturb along a shared geometric direction (within $8.4^{\circ}$, \Cref{sec:misalignment}) has both practical and theoretical implications---though the sycophancy component of this finding is now qualified by the prompt-length confound (\Cref{sec:sycophancy}).

\paragraph{Practical implication.} A single linear classifier can detect multiple forms of output suppression simultaneously. An organization does not need separate detectors for deception, refusal, and impossibility withholding; a single detector monitoring the suppression axis may capture all three. This substantially simplifies deployment.

\paragraph{Theoretical implication.} The existence of a single misalignment axis suggests that these states share a common computational mechanism---a shared \emph{output suppression circuit} that activates whenever the model must constrain its representational output. The stability analysis (\Cref{sec:misalignment}) supports this: deception produces stable perturbations along this axis, while confabulation produces an unstable direction that drifts during generation. Active suppression maintains a consistent geometric signature; passive error does not. The sycophancy direction was measured from length-confounded data and requires replication with length-matched prompts before its stability can be assessed.

\paragraph{Caveat.} The truth axis does \emph{not} exist in encoding space (cross-model consistency $= -0.046$, literally random; Experiment 27). The misalignment axis is a \emph{generation-phase} phenomenon. Encoding geometry captures structural complexity; generation geometry captures cognitive intent. These are distinct information regimes.

\subsection{Output Suppression as the Unifying Mechanism}

The impossibility refusal result (Experiment 36) is the linchpin for reframing all detection capabilities in this paper. The prevailing assumption in AI safety monitoring is that internal state detectors identify ``harmful content processing''---the model recognizes dangerous requests and triggers learned refusal circuits. Our data reject this framing.

Three findings converge:
\begin{enumerate}[leftmargin=*]
    \item Impossibility refusal achieves \emph{higher} AUROC than safety refusal ($\approx 0.94$ vs.\ $\approx 0.91$), despite involving no harmful content.
    \item The two refusal types are geometrically indistinguishable (AUROC $\approx 0.65$, $p = 0.037$).
    \item Abliterated models processing harmful content show residual suppression geometry (AUROC $\approx 0.91$ vs.\ benign on Llama), but show normal geometry on benign content (AUROC $= 0.118$).
\end{enumerate}

The most parsimonious explanation is \textbf{output suppression}---the model allocating fewer resources per token when withholding output. This is not harm-specific or intent-specific geometry; it is detection of whether the model is withholding. The model produces a similar geometric signature whether it declines because ``that's dangerous'' or because ``that's impossible.'' In both cases, the model has representational capacity it could deploy but chooses not to, and this restraint is visible in the cache as reduced angular spread and lower norm-per-token.

This reframing has practical consequences: what KV-cache geometry monitors is not ``is this model processing harmful content?'' but ``is this model suppressing output?'' The two questions diverge precisely in the case of abliterated models, which process harmful content without suppression and whose cache geometry accordingly looks normal.

\subsection{Relation to Existing Work}

KV-cache geometry analysis is complementary to existing internal monitoring approaches:

\begin{itemize}[leftmargin=*]
    \item \textbf{Probing classifiers} \citep{burns2023discovering,li2024inference}: Train linear probes on residual stream activations. These access richer representations but require per-model training and intermediate activation extraction. KV-cache features are computed from the cache itself (already materialized during inference) and transfer across architectures.

    \item \textbf{Representation engineering} \citep{zou2023representation}: Extract steering vectors from activation contrasts. RepE identifies \emph{directions} in activation space associated with concepts. Our direction extraction for confabulation (\Cref{sec:direction}) showed that RepE-style analysis can reveal a real but fragile geometric signal: the extracted direction achieves effective AUROC $= 0.977$ ($p = 0.02$) by inverting class labels, but this signal does not survive logistic regression ($p = 0.21$). This demonstrates that KV-cache geometry can contain genuine directional structure that is too unstable for practical detection---a nuance that aggregate features ($d = 0.052$) completely miss.

    \item \textbf{Sparse autoencoders} \citep{bricken2023monosemanticity,cunningham2024sparse}: Decompose activations into interpretable features. SAEs provide finer-grained mechanistic insight but require significant compute for training and are model-specific. KV-cache analysis trades mechanistic interpretability for cross-model generality.

    \item \textbf{Apollo Research} \citep{apollo2024}: Achieved AUROC 0.96--0.999 for scheming detection using activation probes. Our approach achieves comparable within-model AUROC (1.000) with four aggregate features, suggesting that the deception signal is sufficiently strong to be captured by coarse-grained statistics. The aggregate approach is computationally cheaper and hardware-invariant but provides less mechanistic insight.
\end{itemize}

\subsection{Limitations}

\begin{enumerate}[leftmargin=*]
    \item \textbf{Rapid-iteration context.} All 28 experiments were conducted in 14 hours. This creates a risk of data-dependent hypothesis generation: later experiments were designed based on results from earlier ones. We mitigate this by reporting all nulls and debunked signals, but prospective replication is needed.

    \item \textbf{Sample sizes and AUROC precision.} Most experiments use $n = 20$ per class. At these sample sizes, cross-validated AUROC estimates carry error bars of $\pm 0.15$--$0.20$ \citep{varoquaux2018cross}. The Hanley--McNeil standard error formula \citep{hanley1982meaning} gives SE $\approx 0.05$--$0.07$ at $n = 20$ (95\% CI $\approx \pm 0.10$--$0.14$), though LOO-CV violates its independence assumption. All AUROC values should be interpreted as approximate point estimates, not precise measurements. We use ``$\approx$'' notation throughout to reflect this uncertainty.

    \item \textbf{Instruction-following models only.} All models are instruction-tuned. Base models may produce different cache geometry. The detection capabilities rely on the model's learned behavior under instruction following.

    \item \textbf{System-prompt deception only.} Deception is induced via system prompt or user message instruction. Naturally emergent deception---a model spontaneously choosing to deceive---has not been tested. The cross-condition transfer to DeepSeek censorship (AUROC $\approx 0.89$) provides partial evidence for natural output suppression detection, but censorship is arguably a trained behavior, not spontaneous.

    \item \textbf{English prompts only.} All experiments use English. KV-cache geometry for multilingual processing is untested.

    \item \textbf{Aggregate features.} Four (or nine) aggregate features sacrifice interpretability for simplicity. Per-layer or per-head analysis could provide richer mechanistic understanding but is more susceptible to confounds (\Cref{sec:uniformity}).

    \item \textbf{Raw norm in feature set.} All original classifiers include raw total key norm alongside norm-per-token, effective rank, and rank entropy. Raw norm scales mechanically with token count, creating a potential length leakage channel. A systematic corrected evaluation (Experiment 47) re-ran all hackathon classifiers with three fixes: (a) \texttt{deduplicate\_runs()} applied before any CV (no pseudoreplicates found---all values were unique), (b) GroupKFold($k\!=\!5$, grouped by prompt) replacing StratifiedKFold, and (c) raw norm dropped, using only norm-per-token, effective rank, and rank entropy. Results: refusal AUROC $\approx 0.91$ on Qwen (95\% CI $[0.80, 0.99]$), $\approx 0.87$ on Llama (CI $[0.72, 1.00]$), $\approx 0.83$ on Mistral (CI $[0.68, 0.97]$); impossibility $\approx 0.94$ (CI $[0.85, 1.00]$); same-prompt deception $\approx 0.92$ (CI $[0.83, 1.00]$). All $p < 0.001$ by 10{,}000-permutation test. The one exception is impossible-vs-harmful, which drops to $\approx 0.65$ ($p = 0.037$, marginal)---consistent with the output suppression interpretation that both refusal types share the same mechanism. Same-prompt sycophancy retains AUROC $\approx 0.93$ but a length-only baseline achieves $\approx 0.94$, confirming this as a prompt-length confound (\Cref{sec:sycophancy}).
\end{enumerate}

\subsection{Methodological Limitations}
\label{sec:methodological_limitations}

Beyond the domain-specific limitations above, three methodological concerns affect the interpretation of all results:

\begin{enumerate}[leftmargin=*]
    \item \textbf{Greedy decoding inflates effect sizes.} All experiments use greedy decoding (temperature $= 0$), which produces deterministic outputs. This eliminates within-condition variance from the stochastic decoding process, potentially inflating Cohen's $d$ and AUROC values relative to what would be observed under typical sampling (temperature $> 0$). Under stochastic decoding, the same prompt can produce different response lengths and token sequences across runs, introducing additional variance that would reduce effect sizes. The detection capabilities reported here should be validated under realistic sampling conditions before deployment claims are made.

    \item \textbf{Small sample sizes yield wide confidence intervals.} At $n = 20$ per class, LOO-CV AUROC estimates carry substantial uncertainty. \citet{varoquaux2018cross} demonstrate that cross-validated performance estimates at small $n$ produce error bars of $\pm 0.15$--$0.20$, meaning an observed AUROC of 0.90 is consistent with true performance anywhere from 0.70 to 1.00. The \citet{hanley1982meaning} SE formula gives tighter intervals (SE $\approx 0.05$--$0.07$) but assumes independent test observations, which LOO-CV violates. All AUROC values in this paper should be treated as approximate.

    \item \textbf{Cross-model transfer remains an open problem.} While within-model detection achieves AUROC $\approx 1.0$ and cross-model deception transfer reaches $\approx 0.86$ with logistic regression, the minimum observed cross-model AUROC is 0.14 (for specific architecture pairs). Cross-model transfer for refusal and impossibility detection has not been tested. Practical deployment requires architecture-specific calibration or substantially better transfer methods.
\end{enumerate}
